\paragraph{Theorem F.2.35 (Proved: complete integral homology of a marked gluing).}
For every fully marked F.2 datum there is a terminating exact
integer algorithm computing every integral homology group of
the resulting eight-manifold, including every torsion invariant
factor.  The input consists of the finite homology matrices
$M_1,\ldots,M_k$, free affine numerators $v_i/m_i$ of their
actual orders, the integral clutch vector $\kappa$, and the
chosen primitive cusp marking
\[
 M_0=(M_1\cdots M_k)^{-1},\qquad
 M_0P=P\begin{pmatrix}I_3&I_3\\0&I_3\end{pmatrix},
 \qquad P\in\mathrm{GL}(6,\mathbb Z).                       \tag{F.GC1}
\]
Here $k=2$ or $3$ and the genuine orders are at least two;
order-one sites are first eliminated as in F.2.31.
If a marking is not supplied, the algorithm makes and records
one by Smith reduction.  It does not identify that choice
with an unspecified prescribed marking.
The complete chain prefix through degree nine has at most
4003 generators in any degree.  No spectral-sequence
extension, saturation convention, orientation assumption or
unbounded search enters the homology calculation.

\emph{Proof.}
We supply both global boundary maps which were left outstanding
after F.2.34.  Let $F_k$ be the free group with ordered
generators $s_1,\ldots,s_k$, and let $M(w)$ be the product
specified by a reduced word $w$.  The torus-bundle complement
over the punctured sphere retracts to the corresponding bundle
over a graph.  Its fundamental group and multiplication are
\[
 \Gamma=\mathbb Z^6\rtimes F_k,
 \qquad (x,w)(y,u)=(x+M(w)y,wu).                            \tag{F.GC2}
\]
The fibre and base are aspherical and the base has no
two-cells.  Hence a chosen section identifies the complement
up to homotopy with $B\Gamma$.  These section and peripheral
choices are fixed by the input clutch convention below.

\emph{A finite complement complex with its full bar projection.}
Use reduced words for $F_k$.  The tree resolution of the free
group has $K_0=\mathbb Z[F_k]$ and
$K_1=\bigoplus_i\mathbb Z[F_k]e_i$, with
$de_i=s_i-1$, and no higher terms.  Its explicit augmented
contraction sends a word to its signed path in the Cayley
tree: for each positive letter $s_i$ add its preceding prefix
times $e_i$, and for a negative letter $s_i^{-1}$ subtract
the prefix including that letter times $e_i$.
This is the Fox derivative, and its boundary is $w-1$.
Together with the bar insertion contraction it gives, by
the recursion (F.CA3), a strong contraction of the free-group
bar resolution to this tree resolution.  Every recursion
decreases bar degree and every word is finite.

For a bar simplex $[(x_1,w_1)|\cdots|(x_n,w_n)]$ of $\Gamma$,
put $y_j=M(w_1\cdots w_{j-1})x_j$.
The exact Alexander--Whitney chain map to the torus/free-group
double bar complex is
\[
 \sum_{p+q=n}(-1)^{pq}
 [w_1|\cdots|w_p]\otimes[y_{p+1}|\cdots|y_n].              \tag{F.GC3}
\]
Terms with an identity bar entry are zero.  The horizontal
first face acts on torus vectors by
$\rho(w)=M(w)^{-1}$, and the horizontal boundary has sign
$(-1)^q$.  Notice that $\rho$ is a right action:
$\rho(wu)=\rho(u)\rho(w)$.  This order is essential when
the matrices do not commute.  Formula (F.GC3) follows from
the diagonal bisimplicial bar construction; alternatively,
its chain-map identity follows by pairing adjacent interior
faces, with the first face using
$M(w_1)^{-1}y_j=M(w_2\cdots w_{j-1})x_j$.
The usual shuffle inverse is finite in each degree, so the
map is a chain equivalence.

Apply the torus contraction and (F.CA5), then the free-group
tree contraction and the same finite perturbation formula.
Higher terms decrease group bar degree at least two, whereas
the tree complex only has degrees zero and one.  Thus the
resulting complement complex is
\[
 C_n(E)=\bigwedge^n\mathbb Z^6\ \oplus\
 \bigoplus_{i=1}^k\bigwedge^{n-1}\mathbb Z^6,
 \qquad
 d(e_i\otimes\xi)=(-1)^{n-1}
   (\bigwedge^{n-1}M_i^{-1}-1)\xi.                         \tag{F.GC4}
\]
Its differential alone has this short form; its projection
from group bars retains the finite higher perturbation terms.
Those projection terms are used for every boundary map.

\emph{Finite affine boundary maps.}
For the $i$th finite cap use the orbit groupoid of F.2.33.
Its one-edge homotopy colimit for the affine generator is
the boundary mapping torus.  Vertically contracting it gives
a bounded complex with
$C_n(Y_i)=\bigwedge^n\mathbb Z^6\oplus
\bigwedge^{n-1}\mathbb Z^6$ and differential
$(-1)^{n-1}(\bigwedge^{n-1}M_i-1)$ on the second summand.
Its inclusion into the cyclic bar construction, followed by
the cap projection in (F.CA5), is the full map $Y_i\to B_i$.
In particular this one-edge model avoids an unnecessary
periodic resolution of the seven-dimensional boundary.

Here is an explicit group-bar map for its other inclusion.
Send a torus-groupoid arrow labelled $\lambda$ to the
translation $x^\lambda\in\Gamma$.  With the integral offsets
$z_1(a)$ in (F.CA4), put
$\eta_a=s_i^{-1}x^{z_1(a)}$.
For an arrow $f:a\to b$, the exact identity
\[
 \eta_a\,\psi(gf)=\psi(f)\,\eta_b                        \tag{F.GC5}
\]
is the natural-transformation identity: it is just
$s_i^{-1}x^{M_i\lambda+z_1(b)}
 =x^\lambda s_i^{-1}x^{z_1(b)}$.
For a $q$-simplex, the edge map is $(-1)^q$ times the
alternating prism which inserts $\eta$ at every position,
uses the original arrows before the insertion and the
transformed arrows afterwards.  Its boundary is precisely
the signed difference of the two end maps, by (F.GC5).
Compose this prism with the finite vertical lift of the
one-edge complex, then with (F.GC3) and the complement
projection.  This proves the full chain-map identity.
The offsets record the actual integer affine lift: replacing
$v_i$ by $v_i+m_i\lambda$ changes the offsets and the map,
even if the affine map on the torus has not changed.
An optional unimodular local basis simply applies its matrix
to every translation in (F.GC5).

We make the distinction between a bar cell and the geometric
positive root explicit.  Put
$a_i=z_1(0)=\lfloor v_i/m_i\rfloor$, coordinatewise, and
let $\tau_i$ be the positive one-edge bar cell closed using
the zero-labelled return arrow to object zero.  With the
first-face convention above its geometric loop is
\[
 \tau_i=g_i^{-1}x^{a_i},\qquad
 F_E(\tau_i)=-e_i+M_i^{-1}a_i\in C_1(E),\qquad
 g_i=x^{a_i}\tau_i^{-1}.                                  \tag{F.GC5a}
\]
Thus the bar cell has the opposite peripheral orientation,
with the displayed integer correction when the numerator
has a nonzero integer part.  Formula (F.GC5) sends its
group element to $s_i^{-1}x^{a_i}$; consequently the
geometric positive root in the last expression of (F.GC5a)
maps to $s_i$ exactly.  This is the stated conversion between
the boundary and complement chain conventions.  In a local basis
$Q$, replace $a_i$ in the group and complement formulas by
$Q\lfloor Q^{-1}v_i/m_i\rfloor$.
Writing $\mathcal N_i=\sum_{j=0}^{m_i-1}M_i^j$, the cyclic
degree-two cell has boundary
$m_i\tau_i+w_i-\mathcal N_i a_i$.
Indeed $\tau_i^{m_i}=x^{\mathcal N_i a_i-w_i}$ follows
directly from $g_i^{m_i}=x^{w_i}$.
This independently checks the affine power relation and
the degree-one prism sign, including negative numerators.

There is an optional exact geometric shortcut.  If a verified
local integral basis gives $M_i=\operatorname{diag}(1,D)$ and
$v_i=\epsilon e_0$, $\epsilon\in\{1,-1\}$, eliminate the
first fibre generator from the cap group using
$g^{m_i}=u^\epsilon$.  The cap is the mapping torus with
group $\mathbb Z^5\rtimes_D\mathbb Z$.
The fibre vector $(u,a)$ maps to $x^a g^{\epsilon m_i u}$,
and the root generator maps to $g$.  Substituting this
homomorphism for $\psi$ in the same prism gives the boundary
map to this bounded geometric cap model.  The shortcut is
used only after checking the displayed integral normal form;
the general periodic construction remains available for
every other input.

To check the peripheral convention directly, the universal
cover of $Y_i$ has the root generator $g_i$ with linear part
$M_i$ and affine displacement $v_i/m_i$.  Its integral
translations and $g_i$ map to the same fibre translations
and $s_i$ in (F.GC2).  The full base-circle translation obeys
$g_i^{m_i}=x^{w_i}t$, where $w_i$ is the power vector of
F.2.31.  Filling the disc kills $t$, giving exactly
$s_i^{m_i}=x^{w_i}$.  The straight paths between the chosen
orbit points give (F.GC5), so the prism is this marked map,
including numerators with nonzero integer parts.

\emph{The marked cusp boundary map.}
Write $P=[A\ B]$ in (F.GC1), and set
\[
 w_0=(s_1\cdots s_k)^{-1},\qquad
 c=w_0x^{-\kappa}\in\Gamma.
                                                               \tag{F.GC6}
\]
Send the three angular generators to the columns of $A$,
the three tropical generators to the columns of $B$, and
the base-phase generator to $c$.  Relation (F.GC1) gives
$c b_i c^{-1}=b_i+a_i$ and all the other nilpotent boundary
relations.  Filling kills the angular generators and $c$,
so (F.GC6) gives exactly
$s_1\cdots s_kx^\kappa=1$, in agreement with F.2.31.

This homomorphism has a direct chain implementation on the
full Kuhn diagram.  For an incidence with translation label
$w$, assign the group element $x^{Bw}$.  A vertical phase
$(n,a)\in\mathbb Z\oplus\mathbb Z^3$ maps to $c^nx^{Aa}$.
The incidence shear is $S_w(n,a)=(n,a+n w)$, and
\[
 (c^nx^{Aa})x^{Bw}=x^{Bw}(c^nx^{A(a+n w)}).                \tag{F.GC7}
\]
For a flag of $p$ incidence arrows and a vertical bar
simplex of length $q$, sum over all $(p,q)$ shuffles.
Each incidence step contributes its assigned group element;
each vertical step contributes its phase transformed by
all preceding incidence shears.  Give the shuffle the sign
$(-1)^{\#\{\text{incidence steps before vertical steps}\}}$.
Interior face cancellations use exactly (F.GC7), and adjacent
incidence faces use addition of their translation labels.
This is therefore a chain map from the double bar complex
to $C_*(B\Gamma)$.  Precompose with the finite Kuhn lift of
F.2.34 and postcompose with (F.GC3) and the complement
projection.  The other cusp map is the full logarithmic
collapse already computed in F.2.34.
Both the cusp boundary and the complement are aspherical,
so the displayed marked homomorphism also verifies the
homotopy class of this geometric inclusion.

\emph{The final finite integer calculation.}
Take the mapping cone of the just constructed maps
\[
 \bigoplus_{i=0}^k C_*(Y_i)\longrightarrow
 C_*(E)\oplus\bigoplus_{i=0}^k C_*(B_i),
 \qquad y\longmapsto(f_E(y),-f_B(y)),                     \tag{F.GC8}
\]
where $B_0=W$ is the Kuhn cap.  Gluing collars makes this
the homotopy pushout model of the actual closed manifold.
Consequently its cone computes integral homology directly;
there is no remaining filtration extension to choose.

For clarity the exact finite size bound can be read off
before running any matrix calculation.  Put
$t_n=\sum_{q=0}^{\min(n,6)}\binom6q$, with binomial
coefficients outside their range interpreted as zero.
Let $w_n$ and $y_n$ be the raw cap and boundary ranks of
F.2.34, zero outside their listed ranges.  Then the cone rank
in degree $n$ is
\[
 \binom6n+k\binom6{n-1}+k t_n+w_n+
 k\binom7{n-1}+y_{n-1}.                                 \tag{F.GC9}
\]
For $0\leq n\leq9$ its maximum is 3890 when $k=2$ and
4003 when $k=3$.  Every bar comparison and perturbation
series used to produce these matrices terminates by its
degree filtration.  Word reductions and affine orbit
enumerations are finite, with no imposed size cutoff.
Computing integer kernels followed by Smith normal forms
now gives $H_0,\ldots,H_8$; higher groups vanish by dimension.
Unit cancellations are optional exact integral chain
equivalences which substantially reduce the matrices before
Smith reduction.  This proves the algorithm and the bound.
\hfill$\square$

\emph{Verification and implementation.}
The implementation is in \path{code/fourfolds_global_chains.py}
and \path{code/fourfolds_total_homology.py}.
The generic JSON driver \path{code/fourfolds_marked_gluing.py}
retains the complete supplied marking and affine data, computes
the full group structure of F.2.31, and independently compares
its abelianization with the chain calculation of $H_1$.
The Fox contraction is checked on all 156 explicitly listed
free-group bar simplices through degree three; the exact
Alexander--Whitney and finite-projection identities are
checked on 156 specified semidirect-product simplices with
two noncommuting rank-two matrices.
Four finite boundary comparisons using the independent
seven-dimensional quotient model have acyclic integral
comparison cones.  The more efficient one-edge model is
checked in six stated low-rank cases, including negative,
nonfixed, and nontrivial integer-part numerators; both its
complement comparison and its full Thom cone pass.
The program \path{code/fourfolds_peripheral_signs.py} also
checks (F.GC5a) and the full cyclic degree-two relation in
seven cases, including both actual rank-six caps.  It
reconstructs the positive geometric root as a group element
and checks that its complement bar projection is $+e_i$.
For the Kuhn boundary, every one of its 10048 raw generators
passes the chain-map identity in all degrees zero through
seven.  Its comparison to the corresponding one-generator
complement has an integrally acyclic cone, eliminated by
5088 unit cancellations.  Thus the comparison also checks
the nonzero boundary order-two torsion.
The full eight-dimensional integration for the original
order-$(5,4)$ survivor has passed.  Every marked boundary
column and the complete Mayer--Vietoris differential are
checked.  After 7372 integral unit cancellations, its Smith
homology is $\mathbb Z$ in degrees zero and eight,
$\mathbb Z^4$ in degree four, and zero otherwise, exactly
as in the independent calculation F.2.6.
The corresponding three-point run with logarithmic order
$m=2$ also passes, after 7470 integral unit cancellations.
It gives $H_1=H_6=\mathbb Z/18$,
$H_2=H_5=(\mathbb Z/2)^5$, $H_3=(\mathbb Z/2)^{10}$,
$H_4=\mathbb Z^4\oplus(\mathbb Z/2)^{10}$,
$H_0=H_8=\mathbb Z$ and $H_7=0$, exactly as in the
independent theorem F.2.32.  It reuses only the identical
first-two finite boundary maps: their group bars involve only
the fibre and their own peripheral generator, so adjoining
a third generator changes neither their columns nor their
target row indices.  The third cap, the new cusp word,
the full cone and every Smith group are recomputed.

A third full integration uses
$A=I_3$, $B=\operatorname{diag}(0,1,0)$ in (F.IV2),
actual affine orders $4,4$, numerators $a_1,a_2$, and zero
clutching.  These two linear actions reverse orientation.
Its 7345 integral unit cancellations give free ranks
\[
                  (1,0,4,11,10,1,6,3,0).
\]
The nonzero torsion groups are
\[
\begin{aligned}
 \operatorname{Tor}H_1=\operatorname{Tor}H_5
       &=(\mathbb Z/2)^3\oplus\mathbb Z/4,&
 \operatorname{Tor}H_2&=(\mathbb Z/2)^2,\\
 \operatorname{Tor}H_3=\operatorname{Tor}H_7
       &=\mathbb Z/2,&
 \operatorname{Tor}H_4&=(\mathbb Z/2)^4\oplus(\mathbb Z/4)^5.
\end{aligned}
\]
The independent van Kampen $H_1$ agrees.  Its Euler number
is six, its mod-two Betti vector is
$(1,4,10,14,20,14,10,4,1)$, and its integral top groups
have the required nonorientable form.  Thus this integration
also tests the orientation-reversing part of the stated scope.

The \texttt{--fresh} option bypasses all saved map checkpoints.
The free-bar and semidirect tests, all four periodic-boundary
comparisons, all six one-edge comparison/Thom tests, all seven
peripheral-sign checks, the complete 10048-generator cusp
comparison, and the cached generic survivor computation have
also been independently rerun successfully.  The
reduced-chain, van Kampen $H_1$, Euler and mod-two duality
checks have additionally been independently replayed on
all three complete integration outputs.  The general
algorithm is proved by the explicit chain equivalences and
finite operations above; these finite execution sets are
reported with their exact scope.
